% Authors: Dr. Denys Dutykh (Mathematics Department, Khalifa University of Science and Technology,
% Abu Dhabi, UAE) and Prof. Laurent Vuillon (Univ. Savoie Mont Blanc, CNRS, LAMA, Chambery, France)
\subsection{Relation to earlier work}

The representations studied here belong to an established connection
between Markov numbers, words, and continued fractions. Reutenauer's
Christoffel correspondence parametrizes whole Markov triples
\cite{Reutenauer2009}; the palindromic-closure formula of Reutenauer and
Vuillon gives an occurrence-wise construction of their entries
\cite{ReutenauerVuillon2017}. The underlying central-word and incidence
identities are developed in \cite{BertheDeLucaReutenauer2008}.
These correspondences distinguish occurrences without presupposing that
their largest numerical entries are distinct.

A particularly close arithmetic formulation appears in
\c{C}anak\c{c}{\i} and Schiffler \cite{CanakciSchiffler2020}.
They relate primitive sums of two squares to even palindromic continued
fractions \cite[Theorem~3.10 and Corollary~3.14]{CanakciSchiffler2020}. They
construct the special continued fraction attached to a Markov occurrence
\cite[Theorem~4.1]{CanakciSchiffler2020}, and its right half, through the
proof of \cite[Corollary~4.4]{CanakciSchiffler2020}, is the occurrence pair of
\cref{prop:occurrence-cost} together with its cost $P+Q-2$, as we show after
that proposition.
They conjecture the uniqueness of the representation whose
palindromification is a Markov snake graph
\cite[Conjecture~4.5]{CanakciSchiffler2020} and identify this conjecture
with Markov uniqueness \cite[Theorem~4.8]{CanakciSchiffler2020}. Their
Conjecture~4.6 asks for uniqueness
under the weaker requirement that the half continued fraction use only
the digits $1$ and $2$; they checked both conjectures by computer for all
Markov numbers with denominator below $70$ \cite[Remark~4.7]{CanakciSchiffler2020}.
Our comparison retains every primitive positive representation of the
same integer, with no restriction on its partial quotients. The assertion
is that the occurrence-induced representation minimizes the sum of those
quotients, with the final Euclidean step omitted. It permits ties and
does not assert either of these uniqueness conjectures.

Ordering results compare a different set of objects: the Markov labels
at specified lattice points. Aigner's book records conjectures for rows of
fixed numerator, denominator and sum \cite[Section~10.1]{Aigner2013}; Rabideau and
Schiffler proved the first \cite{RabideauSchiffler2020}, and Lagisquet,
Pelantov\'a, Tavenas, and Vuillon proved all three
\cite{LagisquetEtAl2021}. Their path minimization holds at a fixed lattice
endpoint. Lee, Li, Rabideau, and Schiffler obtain further comparison
inequalities \cite{LeeEtAl2023}, which confine any coincidence of labels to
a narrow cone of slopes; their fixed-row ratio decreases to a constant,
which \cref{sec:row-labels} computes and bounds on both sides from the Cohn
trace, so that each infinite source row lies in a normalized interval. Gaster gives sharp monotonicity ranges using the
stable norm on the modular torus \cite{Gaster2022}. None of these ordering
statements, as stated, minimizes Euclidean cost over the full primitive
two-square fiber of a label. The classical uniqueness theorem for labels
that are prime powers, or twice prime powers, is the case of a fiber with a
single element; \cref{rem:single-representation} derives it from the
results of \cref{sec:occurrences,sec:gaussian}.

Two of our submitted articles treat neighboring questions. The first
studies the same Euclidean cost on the representations produced by a single
Gaussian reflection \cite{DutykhVuillon2026Endpoint}, and contains the cost
identity of \cref{prop:occurrence-cost} in its source dictionary; the second
studies palindromes in the Cohn monoid \cite{DutykhVuillon2026Reflections},
and contains \cref{lem:fixed-cost-injectivity} as the case of equal sums in
the proof of its proposition on the lengths of directing words. Their full texts can be consulted in the repository of
this series, given at the end of the article.

Related extremal problems fix the length and sum of the entries of a
continuant; Kan determines the numerical maximum under these constraints
\cite[Theorem~5]{Kan2012}.\footnote{For sequences of length $t$ and sum $S$,
the published statement assumes $2\leqslant t\leqslant S<nt$, where $n$, a bound
on the entries, belongs to the next theorem of that article; the corrected
version \cite[Theorem~3.5]{Kan2021} states the result for
$2\leqslant t\leqslant S$.}
Our proof derives the required path bounds
directly. In its infinite part, a parity order for the full normalized
ratio reduces whole boxes of gap lengths to exact corner comparisons.
The finitely many remaining labels are settled through their complete
Gaussian factorizations, so that an infinite comparison argument is
joined to a finite piece of exact arithmetic. The resulting thirteen-row
theorem assumes no upper bound on the denominator or on the Markov
value. Its injectivity consequence applies when both occurrences
belong to the proved rows; it does not exclude a collision with an
occurrence outside them.
